\documentclass[tikz,border=10mm]{standalone}
\usetikzlibrary{arrows.meta, calc}

% Поддержка русского языка для компилятора pdfLaTeX
\usepackage[T2A]{fontenc}
\usepackage[utf8]{inputenc}
\usepackage[russian]{babel}

\begin{document}
\begin{tikzpicture}[
    patch/.style={rectangle, draw=blue!80!black, fill=blue!10, minimum width=10mm, minimum height=8mm, font=\small\bfseries},
    divider/.style={circle, draw=red!80!black, fill=red!20, inner sep=1.5pt, font=\tiny},
    line50/.style={draw=black, line width=1.5pt},
    line70/.style={draw=black!70, line width=1.0pt}
]

    % Заголовок
    \node[font=\Large\bfseries] at (0, 7.5) {Топология СВЧ-распределительной сети H-Tree (4x4)};

    % 16 Патч-антенн с безопасным именованием через скобки
    \foreach \x in {-4.5, -1.5, 1.5, 4.5} {
        \foreach \y in {-4.5, -1.5, 1.5, 4.5} {
            \node (P_{\x}_{\y}) at (\x,\y) [patch] {патч};
        }
    }

    % ЦЕНТРАЛЬНЫЙ ВХОД (Уровень 1)
    \node (IN) at (0, -7) [below, font=\small\bfseries] {Вход (50 $\Omega$)};
    \draw[line50, -{Latex[scale=1.2]}] (0, -6.5) -- (0, 0);
    \node (D1) at (0,0) [divider] {D1};

    % УРОВЕНЬ 2
    \node (D2-L) at (-3, 0) [divider] {D2};
    \node (D2-R) at (3, 0) [divider] {D2};
    \draw[line50] (D1) -- (D2-L);
    \draw[line50] (D1) -- (D2-R);

    % УРОВЕНЬ 3
    \node (D3-TL) at (-3, 3) [divider] {D3};
    \node (D3-BL) at (-3, -3) [divider] {D3};
    \draw[line50] (D2-L) -- (D3-TL);
    \draw[line50] (D2-L) -- (D3-BL);
    
    \node (D3-TR) at (3, 3) [divider] {D3};
    \node (D3-BR) at (3, -3) [divider] {D3};
    \draw[line50] (D2-R) -- (D3-TR);
    \draw[line50] (D2-R) -- (D3-BR);

    % Квадрант Top-Left
    \node (D4-TL1) at (-4.5, 3) [divider] {D4};
    \node (D4-TL2) at (-1.5, 3) [divider] {D4};
    \draw[line70] (D3-TL) -- (D4-TL1); \draw[line70] (D3-TL) -- (D4-TL2);
    \draw[line70] (D4-TL1) -- (P_{-4.5}_{4.5}); \draw[line70] (D4-TL1) -- (P_{-4.5}_{1.5});
    \draw[line70] (D4-TL2) -- (P_{-1.5}_{4.5}); \draw[line70] (D4-TL2) -- (P_{-1.5}_{1.5});

    % Квадрант Bottom-Left
    \node (D4-BL1) at (-4.5, -3) [divider] {D4};
    \node (D4-BL2) at (-1.5, -3) [divider] {D4};
    \draw[line70] (D3-BL) -- (D4-BL1); \draw[line70] (D3-BL) -- (D4-BL2);
    \draw[line70] (D4-BL1) -- (P_{-4.5}_{-1.5}); \draw[line70] (D4-BL1) -- (P_{-4.5}_{-4.5});
    \draw[line70] (D4-BL2) -- (P_{-1.5}_{-1.5}); \draw[line70] (D4-BL2) -- (P_{-1.5}_{-4.5});

    % Квадрант Top-Right
    \node (D4-TR1) at (1.5, 3) [divider] {D4};
    \node (D4-TR2) at (4.5, 3) [divider] {D4};
    \draw[line70] (D3-TR) -- (D4-TR1); \draw[line70] (D3-TR) -- (D4-TR2);
    \draw[line70] (D4-TR1) -- (P_{1.5}_{4.5}); \draw[line70] (D4-TR1) -- (P_{1.5}_{1.5});
    \draw[line70] (D4-TR2) -- (P_{4.5}_{4.5}); \draw[line70] (D4-TR2) -- (P_{4.5}_{1.5});

    % Квадрант Bottom-Right
    \node (D4-BR1) at (1.5, -3) [divider] {D4};
    \node (D4-BR2) at (4.5, -3) [divider] {D4};
    \draw[line70] (D3-BR) -- (D4-BR1); \draw[line70] (D3-BR) -- (D4-BR2);
    \draw[line70] (D4-BR1) -- (P_{1.5}_{-1.5}); \draw[line70] (D4-BR1) -- (P_{1.5}_{-4.5});
    \draw[line70] (D4-BR2) -- (P_{4.5}_{-1.5}); \draw[line70] (D4-BR2) -- (P_{4.5}_{-4.5});

    % Легенда
    \matrix [draw, below left, fill=white, opacity=0.9] at (7,-5.5) {
      \node [patch, scale=0.5, label=right:Излучатель (патч)] {}; \\
      \node [divider, scale=0.7, label=right:Делитель (D1-D4)] {}; \\
      \node [draw=black, line width=1.5pt, minimum width=6mm, label=right:Основная линия 50 Ohm] {}; \\
      \node [draw=black!70, line width=1.0pt, minimum width=6mm, label=right:Преобразователь импеданса] {}; \\
    };

\end{tikzpicture}
\end{document}
